\section{The resource-transition benchmark in full}\label{si:benchmark}

Main-text Section~\ref{main-sec:benchmark} summarizes this synthetic illustration. Its biomass component is a Schaefer equation in model biomass units (MBU); the income-margin and reserve coordinates are normalized paths stipulated by the witness, not outputs of an economic model. The model is not fitted to a fishery. Published Northern-cod values are used only for the optional historical dimensional map below.

\paragraph{State, scales, and disturbance} Let $B_m(t)$ denote spawning biomass in MBU, with $t$ measured in years and $\varsigma(B_m)=rB_m(1-B_m/K)$, where $r=4\,\mathrm{yr}^{-1}$, $K=10$ MBU, and $B_{\mathrm{lim}}=2$ MBU. The ecological margin is $s_1=(B_m-B_{\mathrm{lim}})/(1\,\mathrm{MBU})$. The second coordinate is a normalized income margin $s_2=\Delta\Pi/\Pi_0$, for a declared scale $\Pi_0>0$; no price, cost, or revenue equation is supplied. The reserve $x$ is in normalized fund units, with staged cost $c=1$. A heatwave strike, when present, subtracts $\delta_0=1/2$ MBU from biomass at $t=1/2$; it does not affect $s_2$. The weighted dashboard acts on the dimensionless normalized coordinates, not on raw tonnes and currency.

\paragraph{The exact first-half schedule and the tube envelope} At $(x,s_1,s_2)=(1/2,6/5,6/5)$, $B_m(0)=16/5$ MBU. On $0\le t\le1/2$, set $B_{\mathrm{line}}(t)=16/5-3t$ and
\[
H_m(t)=3+\varsigma(B_{\mathrm{line}}(t)),\qquad \dot B_m=\varsigma(B_m)-H_m(t).
\]
The benign first half follows $B_{\mathrm{line}}$ exactly. The quota decreases from $1463/125=11.704$ to $2161/250=8.644$ MBU/yr, under the stipulated $H_{\max}=13$ MBU/yr. At the midpoint the benign biomass is $17/10$ MBU; the heatwave branch drops to $6/5$ MBU. The certified floor paths are $s_1:6/5\to-3/10\to6/5$ without the event and $s_1:6/5\to-4/5\to6/5$ with it. The \act{slow} income-margin excursion is prescribed analogously, with biomass held at $16/5$ MBU. The \act{staged} path spends one fund unit and raises each normalized margin by $1/4$; its income and reserve trajectories are stipulated, not derived from cash-flow equations. The abstract action tubes and destination rules remain those listed in main-text Table~\ref{main-tab:actions}.

With zero harvest after the midpoint, biomass follows the logistic equation rather than the piecewise-linear recovery envelope. From midpoint value $b$, the endpoint is
\[
B_m(1)=\frac{10}{1+(10/b-1)e^{-2}}.
\]
This gives $6.021359\ldots$ MBU from $b=1.7$ and $5.018924\ldots$ MBU from $b=1.2$, not the abstract successor $16/5$ MBU. The straight-line recovery is only a conservative lower envelope: until the realized stock reaches $5$ MBU, $\varsigma(B_m)\ge\varsigma(6/5)=528/125>4$ MBU/yr, exceeding the envelope's largest recovery rate; afterward the realized stock remains above $5$ while the envelope is at most $16/5$. To use the abstract successor in a physical model, an additional end-of-interval control or reset must be specified. The comparison proves a biomass lower envelope, not the full two-coordinate transition or terminal attainability.

\paragraph{Income-margin dynamics and disturbance treatment} For the displayed witness, the illustrative schedules are: \act{no-switch} holds both normalized margins; \act{fast} holds $s_2$ while $s_1$ follows its pulse-and-recovery tube; \act{slow} holds $s_1$ while $s_2$ follows the analogous depth-2 excursion; and \act{staged} increases both margins linearly by $1/4$ while the fund is spent. The heatwave affects only $s_1$, and the declared no-event branch removes that biomass strike. These are complete specifications of the benchmark's second coordinate at the level of the witness tube, but not an empirically derived income model.

\paragraph{Dashboard verdict} On the heatwave branch of \act{fast}, the normalized sum is $(12/5,2/5,12/5)$, so the equal-weight margin remains positive while the biomass floor is breached. At the canonical point, \act{fast} is licensed for $r=w_2/w_1\ge2/3$ and \act{slow} for $r\le3/2$; the intervals overlap, and a biomass-only weight rejects \act{fast} but can license \act{slow}. The claim is weight-adaptive state acceptance, not acceptance of one fixed trajectory under every weight. A biomass-limit breach is not, by itself, a claim of biological collapse. For states outside $\Vtyp$ with $x<1$, the reserve shortfall is $\kappa^*=1-x$ under the witness's staged action.

\begin{table}[!htbp]
\centering\small
\caption{The witness coordinates read as quantities in the illustrative benchmark.}
\label{tab:fisheries}
\begin{tabularx}{\textwidth}{@{}>{\raggedright\arraybackslash}p{0.3\textwidth}X@{}}
\toprule
Object & Illustrative reading \\
\midrule
Typed margins $s_i\ge0$ & normalized biomass margin $s_1$ and separately normalized income-margin coordinate $s_2$ \\
Scalarized operator $E_w$ & dashboard $w\cdot(s_1,s_2)$ applied to declared dimensionless margins \\
Licensing thresholds $\rho_1,\rho_2$ & weight-ratio ranges licensing the pulse (\act{fast}) and gradual (\act{slow}) plans \\
Acceptance gap $\FP$ & dashboard-certified yet typed-inadmissible transitions under the declared menu and disturbances \\
Rescue set $R$; threshold $\kappa^*$ & states served by the staged plan and normalized fund increment needed to reach its cost \\
Blend versus concatenation & fractional physical control versus the separately stipulated union of complete primitive tubes \\
Exact-tube semantics & checking all states in the declared tube, including the midpoint event \\
CES elasticity $\sigma^*$ & threshold in the constructed normalized-margin family, not a fitted biological parameter \\
\bottomrule
\end{tabularx}
\end{table}

\paragraph{Northern-cod anchor and dimensional map} DFO Science Advisory Report 2016/026, Table A2 (PDF p.~16), reports the assessment-vintage Northern-cod LRP as $884.6$ kt, the rounded 1983--1989 mean spawning-stock biomass of that series ($884.58$ kt from a total of $6192.06$ kt). The 2015 value is $298.65$ kt, about $33.8\%$ of that historical LRP. This is not asserted to be a current-2026 threshold. If, only as a unit comparison, $2$ MBU is mapped to $884.6$ kt, then $1$ MBU equals $442.3$ kt:
\begin{center}\small
\begin{tabular}{@{}lrr@{}}\toprule
Model quantity & MBU value & mapped physical value \\
\midrule
Limit $B_{\mathrm{lim}}$ & $2$ & $884.6$ kt \\
Opening biomass $B_m(0)$ & $16/5$ & $1415.36$ kt \\
Adverse trough & $6/5$ & $530.76$ kt \\
Carrying capacity $K$ & $10$ & $4423$ kt \\
Heatwave strike $\delta_0$ & $1/2$ & $221.15$ kt \\
Capacity $H_{\max}$ & $13$ MBU/yr & $5749.9$ kt/yr \\
\bottomrule\end{tabular}
\end{center}
These large implied capacities make clear that the map is dimensional only; it does not make the model a plausible or fitted cod fishery. It is also important to preserve the index normalization: $B_m/B_{\mathrm{lim}}=1+s_1/2$, while the CES construction uses $\lambda_1=1+s_1$. If a physical biomass margin is $s_{1,\mathrm{phys}}=442.3s_1$ kt, then $\lambda_1=1+s_{1,\mathrm{phys}}/442.3$, not $1+s_{1,\mathrm{phys}}$. The income coordinate uses its own scale, and the reserve has no physical currency unit here. No model parameter is estimated from DFO data, and the witness makes no empirical prediction.

\paragraph{Institutional interpretation of the reserve} The normalized threshold $\kappa^*=1-x$ can be read as an abstract adjustment-fund shortfall. Actual buy-backs, transition assistance, or quota banking require a monetary accounting model and an implementation mechanism not supplied by this illustration; these are examples of possible interpretations, not consequences of the theorem.

\section{Practitioner material}\label{si:practice}

\subsection{Terminology map}

\begin{table}[!htbp]
\centering\small
\caption{Terminology map: assessment objects across three neighbouring vocabularies (1/2). The entries are correspondences of mathematical role, not analogies.}
\label{tab:translation}
\begin{tabularx}{\textwidth}{@{}>{\raggedright\arraybackslash}p{0.19\textwidth}>{\raggedright\arraybackslash}X>{\raggedright\arraybackslash}X>{\raggedright\arraybackslash}X@{}}
\toprule
 & MCDM / operations research & Control / viability theory & Environmental modelling / governance \\
\midrule
Typed floors $s_i\ge0$, path-wise & coordinate-wise criteria; per-criterion lower bounds along the path & state constraints; safe set; tube constraints & planetary boundaries; safe operating spaces; regulatory floors (for example, a biomass LRP) \\
\addlinespace
Scalarized operator $E_w$ & weighted-sum acceptance; weight space $\Wp$ & scalarized robust constraint on an aggregate output & composite sustainability index; the dashboard \\
\addlinespace
Licensing thresholds $\rho_1,\rho_2$ & weight intervals of acceptance; robustness of a recommendation to weights & parameter-dependent feasibility margins & indicator weighting schemes; stakeholder weight elicitation \\
\addlinespace
Acceptance gap $\FP$ & every weight licenses some plan, but no single plan meets all typed floors & weight-wise feasibility with no common admissible plan & aggregate acceptance under the stated rule despite typed rejection (menu-relative) \\
\bottomrule
\end{tabularx}
\end{table}

\begin{table}[!htbp]
\centering\small
\caption{Terminology map, continued (2/2).}
\label{tab:translation-controls}
\begin{tabularx}{\textwidth}{@{}>{\raggedright\arraybackslash}p{0.19\textwidth}>{\raggedright\arraybackslash}X>{\raggedright\arraybackslash}X>{\raggedright\arraybackslash}X@{}}
\toprule
 & MCDM / operations research & Control / viability theory & Environmental modelling / governance \\
\midrule
Rescue set $R$; threshold $\kappa^*$ & alternatives repairable by a budgeted side payment & reachability under an auxiliary resource input & transition finance: adjustment funds or fleet buy-backs \\
\addlinespace
Impossibility region $I$ & accepted for every weighting, yet no common plan is feasible & empty common safe-action set; kernel exclusion & aggregate acceptance with a floor crossed mid-transition \\
\addlinespace
Fractional blend vs.\ complete-excursion concatenation & physical fractional controls vs.\ the stipulated union of full primitive tubes & relaxed control vs.\ a declared switching class & blended schedule vs.\ complete primitive excursions \\
\addlinespace
Exact-tube semantics & robust path-wise feasibility over a scenario set & viability tubes; worst-case reachability & avoiding transient breaches during a transition \\
\bottomrule
\end{tabularx}
\end{table}

\subsection{A ten-point audit checklist}

For an agency using a composite dashboard to certify a sustainability transition, the four design tests of main-text Section~\ref{main-sec:design} operationalise as ten audit points.
\begin{enumerate}[leftmargin=2em]
\item \textbf{Object of certification:} state what is being certified -- a static state, an endpoint, or the existence of a path-wise safe plan.
\item \textbf{Binding floors:} list the floors that are legally binding, irreversible, or lack verified substitutability; these require noncompensatory (Protocol~4) treatment.
\item \textbf{Policy quantifier:} declare whether the approval rule assumes one robust plan or lets the plan adapt to the weighting.
\item \textbf{Weight-dependent plans:} if weight-adaptive, report the plan associated with each relevant weight range (the thresholds $\rho_1,\rho_2$ are the pattern).
\item \textbf{Path-wise evaluation:} evaluate the floors along the full transition, not only at assessment snapshots.
\item \textbf{Disturbance structure:} specify whether shocks are common to all plans, action-indexed, or coupled across floors (Table~\ref{tab:disturbance}); the separation's scope and the remedy ride on the class.
\item \textbf{Menu completeness:} acknowledge that negative feasibility results are menu-relative unless the action space is demonstrably exhaustive.
\item \textbf{Blended policies:} do not infer convexification from mathematical convenience; verify that fractional effort or quota allocation is administratively implementable.
\item \textbf{Rescue margins:} if a resource-controlled action exists, report the minimum increment $\kappa^*$ that converts impossibility into rescue, and whether it is finite under the declared disturbance class.
\item \textbf{Mandatory reporting:} require that component trajectories and their minimum margins, the reserve and the acceptance protocol be reported alongside the aggregate, so that the typed verdict can be independently reconstructed; floors alone do not fix the side of the rescue threshold.
\end{enumerate}

\subsection{A minimum report}

A minimum report for a transition assessment should include: (1) typed floors and their provenance; (2) aggregate weights or the weight family; (3) policy quantifiers; (4) the disturbance set and class; (5) the observation and implementation model; (6) exact versus conservative tube status; (7) terminal maintainability; (8) action-exhaustion status; (9) uncertainty and approximation bounds; (10) normative premises.

\subsection{Data requirements for empirical application}

Empirical application requires, in addition to the typed floors, disturbance set, plan menu, tube model and destination-maintainability witness, five further specifications.
\begin{enumerate}[leftmargin=2em]
\item \textbf{Action-set completeness status.} A negative certificate over a finite menu proves impossibility only relative to that menu. The application must state whether $\mathcal A$ is exhaustive, a listed policy menu, a sampled subset, or an inner approximation. If incomplete, the verdict is ``no safe transition exists among the listed plans'', not ``no safe transition exists''.
\item \textbf{Calibration and identifiability data.} Parameter estimates, uncertainty sets, structural alternatives, validation data, observation error, model discrepancy, missing-data treatment, disturbance dependence. A specified disturbance set is necessary but not sufficient; it must be justified, or the verdict marked conditional on model credibility.
\item \textbf{Policy and authority data.} Who may select each plan, what information they possess, decision timing, enforcement assumptions, compliance uncertainty, strategic responses, resource and legitimacy constraints.
\item \textbf{Threshold-uncertainty semantics.} A floor is rarely known exactly. Distinguish the deterministic floor $s_i\ge0$ from robust threshold safety $s_i(z;\vartheta)\ge0$ for all $\vartheta\in\Theta$, and from probabilistic or confidence-level versions if admitted.
\item \textbf{Destination maintainability status.} Whether the destination-maintainability witness is established physically, by simulation, or by assumption.
\end{enumerate}
As a dimensional illustration only, the benchmark optionally maps $2$ MBU to the historical 2016 Northern-cod LRP of $884.6$ kt; this is not a parameter calibration or a claim about the current threshold.

\subsection{Endpoint evaluation on the witness}

The typed-endpoint operator $E_{\mathrm{end,typ}}(z)=\{a:\ \forall d,\ \End(a,d)\subseteq S\text{ and }\Succ(a,d)\subseteq G\}$ satisfies $\Etyp\subseteq E_{\mathrm{end,typ}}\subseteq\Eend$. On the witness, \act{fast}'s endpoint values of $s_1$ equal the initial $s_1$ and its successor lies in $G$ whenever $x\ge0$, so $E_{\mathrm{end,typ}}(z)\ne\varnothing$ at every $z\in X_0$, while $\Etyp(z)\ne\varnothing$ holds exactly on $\{x\ge1\}\cup\{s_1\ge2\}\cup\{s_2\ge2\}$: on the stricken slice $\Q\cap\{x<1\}$ the endpoint evaluation is nonempty where the typed operator is empty -- the photograph admits what the trajectory forbids. Reporting regimes that evaluate endpoints only (audited annual snapshots) therefore license transitions that violate typed floors mid-interval without detection; no empirical claim about any particular reporting regime is made.

\begin{figure}[!htbp]
\centering
\includegraphics[width=0.9\linewidth]{figS_weight_intervals.png}
\caption{Weight-interval view at the witness point $(x,s_1,s_2)=(\tfrac12,\tfrac65,\tfrac65)$. \act{fast} is licensed exactly for $r\ge\rho_1=\tfrac23$ and \act{slow} exactly for $r\le\rho_2=\tfrac32$. The two intervals together cover every weight ratio $r\ge0$, but neither covers it alone. On the rescue set, \act{staged} supplies the full interval and closes the gap.}
\end{figure}
